\documentclass[10pt]{article}
\usepackage[a5paper,margin=13mm]{geometry}
\usepackage{fontspec}
\usepackage{xeCJK}
\usepackage{xcolor}
\usepackage{booktabs}
\usepackage{tabularx}
\usepackage{enumitem}
\usepackage{hyperref}
\setmainfont{TeX Gyre Pagella}
\setsansfont{TeX Gyre Heros}
\setCJKmainfont{Noto Serif CJK TC}
\setCJKsansfont{Noto Sans CJK TC}
\definecolor{ink}{HTML}{17231F}
\definecolor{green}{HTML}{2B6B57}
\definecolor{orange}{HTML}{D96B3B}
\hypersetup{colorlinks=true,urlcolor=green}
\setlength{\parindent}{0pt}
\setlength{\parskip}{5pt}
\setlist{nosep,leftmargin=*}
\pagestyle{plain}
\begin{document}

{\sffamily\small\color{green} BILINGUAL LECTURE PACK · PROJECT-OWNED SAMPLE}

{\LARGE\bfseries\color{ink} Why provenance matters in a\\multilingual knowledge graph}

{\small English → Traditional Chinese · 00:41.191 · synthetic source}

\section*{Summary · 摘要}
A useful knowledge-graph relation keeps its evidence close: the claim, the exact supporting passage, and the source-file fingerprint. In multilingual work, the original phrase, translation, and uncertainty should remain visible together so a reader can inspect and correct the relation.

一條有用的知識圖譜關係，應把證據留在旁邊：主張、支持它的確切原文段落，以及來源檔案指紋。處理多語資料時，原文、譯文和不確定之處也應並列保留，讓讀者能查驗和修正關係。

\section*{Five key ideas · 五個重點}
\begin{enumerate}
  \item A relationship is only as useful as its evidence. \hfill 關係的價值取決於證據。
  \item Preserve the exact supporting passage. \hfill 保留支持主張的確切原文。
  \item Fingerprint the source file. \hfill 記錄來源檔案指紋。
  \item Keep translation uncertainty visible. \hfill 讓翻譯的不確定之處可見。
  \item Provenance makes correction possible. \hfill 來源記錄讓查驗與修正成為可能。
\end{enumerate}

\section*{Aligned terms · 對照詞彙}
\small
\begin{tabularx}{\textwidth}{@{}lXr@{}}
\toprule
English & Traditional Chinese & First use \\
\midrule
knowledge graph & 知識圖譜 & 00:00 \\
evidence & 證據 & 00:00 \\
claim & 主張 & 00:09 \\
exact passage & 確切原文段落 & 00:09 \\
source-file fingerprint & 來源檔案指紋 & 00:09 \\
translation & 翻譯／譯文 & 00:15 \\
uncertainty & 不確定之處 & 00:17 \\
provenance & 資料來源記錄 & 00:35 \\
inspectable & 可供查驗 & 00:35 \\
\bottomrule
\end{tabularx}

\newpage
\section*{Timestamp trail · 時間索引}
\begin{itemize}
  \item \textbf{00:00--00:05} — Evidence gives a graph relation its value.
  \item \textbf{00:09--00:15} — Keep claim, passage, and source fingerprint together.
  \item \textbf{00:17--00:25} — Show original, translation, and uncertainty side by side.
  \item \textbf{00:27--00:34} — A reader follows the link back and checks the wording.
  \item \textbf{00:35--00:40} — Provenance supports inspection, not automatic truth.
\end{itemize}

\section*{Review questions · 複習問題}
\begin{enumerate}
  \item Which three items should stay together in the graph?\\
  圖譜中的哪三樣東西應放在一起？
  \item Why should translation uncertainty remain visible?\\
  為什麼應保留翻譯的不確定之處？
  \item What can provenance establish, and what can it not establish?\\
  來源記錄能證明什麼，又不能證明什麼？
\end{enumerate}

\section*{Review and provenance}
\small
The English transcript was checked against the included project-owned script and source audio. The Traditional Chinese is an assisted translation with terminology and alignment review; it is not certified and has not received independent native-language review. Translation decisions remain in \texttt{aligned-source.json} and \texttt{issue-ledger.md}.

Source and sample content: LazyingArt LLC, CC BY 4.0. This is process evidence, not customer work. File hashes and build versions are recorded in the delivery manifest.

\end{document}
